\section{Conclusions}
\label{sec:conclusions}
% - works? really works?

The \textit{Kitsurai} project was developed with the goal of creating a distributed key-value store
capable of offering predictable and guaranteed performance in multi-tenant environments,
a requirement often overlooked in major existing open-source solutions.

The developed system has ample margin for improvements, but demonstrates how the core techniques
can be used to create a database fit for real-time guarantees.

The experimental results confirm the validity of the approach, as the system:
\begin{itemize}
    \item managed to maintain controlled performance even under high load and in the presence of noisy clients,
    ensuring effective isolation between tenants and guaranteeing each a minimum share of bandwidth and an acceptable maximum latency;
    \item managed to scale to within variance of the expected rate for multiple configurations.
\end{itemize}
Even thought testing was conducted on non-enterprise hardware due to availability,
we could expect at least to match the result here presented when switching to high-end enterprise hardware.

\subsection{Limitations and Future Work}
\label{subsec:limitations+future-work}

Despite the promising results, the current implementation has some limitations that provide opportunities for future work.

The developed algorithm has been developed for a static set of peers,
but dynamic peers are a possibility.
In particular, a catch-up protocol for reintegrating nodes after prolonged outages
or after joining the cluster would have to be developed.

On the same note, table's bandwidth parameter are static and cannot be changed.
An extension to the system could revise the table allocation and key distribution to allow changes.
% A possibility could be to allow nesting of the allocation rings...
Relaxing the uniform key access assumption, a system like Dynamo's split-for-hot could then be implemented as an internal auto-scaler,
by using the developed dynamic table API\@.

Encryption is already supported by the underlying HTTP/2 server's implementation,
exposing the configuration to the CLI is most of the work to enable TLS on both the client's HTTP and for the internal RPC server.
Further testing could also be done to test HTTP/3 and QUIC as an alternative to HTTP/2, % TODO: Cita?
which was not done due to the requirement of TLS\@.

Strategies for optimal capacity management of nodes, especially when tables have unbalanced requirements in terms of bandwidth or storage space,
coupled with the likely need to improve the bandwidth limitation model to allow separate limits for get and set operations.
An approach like that of X-Stor's RUs or DynamoDB's RCUs and WCUs could be a nice addition to the current capabilities.

Many other features can interestingly be built on top of the existing system, as a series of middlewares:
\begin{itemize}
    \item An access control model for fine-grained permission management by using industry standard technologies like OAuth/JWT; % TODO: Cita qualcosa
    \item Built-in consistency options, like timestamps or Dynamo's vector clocks, could also be implemented transparently;
    \item Assigning custom names to tables could also be built on top, storing the table's translations in the database itself.
\end{itemize}
